\section{Применение параллелизма к верификации нейронных сетей}
\label{sec:parallel-verifiers}

В данном разделе обосновывается выбор Tensor Parallelism и FSDP для
задачи верификации, анализируется их применимость к алгоритмам bound
propagation и описываются трудности интеграции с полной верификацией.

\subsection{Почему Tensor Parallelism подходит для верификации}

Алгоритмы распространения границ (CROWN, $\alpha$-CROWN) сводят
вычисление bounds к цепочке матричных умножений, выполняемых обратным
проходом по вычислительному графу сети. Для линейного слоя с весовой
матрицей $W$ обратный проход CROWN выполняет умножение
$A \leftarrow A \cdot W$, где $A$~--- матрица коэффициентов линейной
релаксации. Эта операция структурно идентична backward pass при обучении.

Из этого следуют три обстоятельства, делающие TP естественным выбором.
Во-первых, для полносвязных сетей весовые матрицы составляют основную
часть потребления GPU-памяти, и TP разрезает именно их, обеспечивая
близкое к линейному снижение с ростом числа GPU.
Во-вторых, коммуникационный паттерн совпадает с обучением: в CROWN
backward каждый слой вносит один \texttt{AllReduce}, все промежуточные
вычисления выполняются локально.
В-третьих, TP сохраняет единый вычислительный граф на всех GPU~---
в отличие от Pipeline Parallelism, который разрезает граф на стадии.
Это критично для \texttt{auto\_LiRPA}, строящей граф bound propagation
поверх исходного.

\subsection{Влияние Tensor Parallelism на точность bounds}
\label{sec:tp-precision-loss}

Ключевой недостаток TP~--- потеря точности bounds, растущая с числом
шардированных зон. Эта потеря не является ошибкой реализации: она
фундаментально обусловлена тем, как bound propagation взаимодействует
с частичными весовыми матрицами.

\subsubsection{Обратное распространение в CROWN}

CROWN вычисляет линейные bounds выхода сети как функцию входа.
Для сети с $d$ слоями и весовыми матрицами $W_1, \ldots, W_d$ обратный
проход строит A-матрицу~--- произведение линеаризованных весов:
$$A = \Lambda_d W_d \cdot \Lambda_{d-1} W_{d-1} \cdots \Lambda_1 W_1,$$
где $\Lambda_i = \mathrm{diag}(\lambda_i)$~--- диагональная матрица
коэффициентов линейной релаксации ReLU (значения из $[0, 1]$,
определяемые по промежуточным bounds нейронов).

Итоговые bounds принимают вид:
$$\ell \leq A \cdot \mathbf{x} + \mathbf{b} \leq u,$$
где $\mathbf{x}$~--- входной вектор в пределах $\varepsilon$-окрестности.
Ключевое свойство: $A$~--- единая линейная функция от \textit{исходного}
входа, сохраняющая зависимости между переменными на протяжении всех слоёв.

\subsubsection{Проблема CROWN в шардированных зонах}

При TP каждая пара Column--Row формирует <<шардированную зону>>~---
подграф, в котором активации имеют размерность $h/P$ на каждом GPU
вместо~$h$.

Пусть слой $W_i \in \mathbb{R}^{h \times h}$ шардирован по столбцам:
GPU~$r$ хранит $W_i^{(r)} \in \mathbb{R}^{(h/P) \times h}$.
При CROWN backward A-матрица предыдущего слоя имеет полную
размерность~$h$:
$$A \leftarrow A \cdot \Lambda_i W_i.$$
Для этого умножения каждому GPU необходима \textit{полная} матрица $W_i$,
однако каждый GPU хранит лишь $W_i^{(r)}$~--- $h/P$ строк.
Выполнить корректное умножение без коммуникации невозможно.

Для пары Column--Row без промежуточных слоёв проблема решаема:
Row-слой выполняет \texttt{AllReduce}, восстанавливая полный результат.
Но если между Column и Row находятся ReLU (как в типичных сетях
глубже 2 слоёв), ситуация усложняется: для выбора коэффициентов
$\lambda_i$ CROWN необходимо знать \textit{промежуточные bounds}
нейронов в шардированной зоне, а их точное вычисление потребовало бы
отдельного обратного прохода с коммуникацией на каждом шаге~---
что сводит на нет экономию памяти.
Единственная практичная альтернатива~--- использовать IBP для
промежуточных bounds внутри зон, поскольку IBP работает послойно
с локальными весами без A-матриц.

\subsubsection{Почему IBP существенно менее точен}

IBP и CROWN вычисляют bounds для одной и той же задачи, но
принципиально различаются в обращении с зависимостями между
переменными.

\textbf{CROWN} поддерживает глобальную линейную зависимость
$y = A \mathbf{x} + \mathbf{b}$, где $A$ учитывает влияние каждого
входного компонента через \textit{все} промежуточные слои.
Bounds вычисляются однократно:
$$lb = \sum_j \min(A_j x_j^L, A_j x_j^U) + b, \quad
  ub = \sum_j \max(A_j x_j^L, A_j x_j^U) + b.$$

\textbf{IBP} вычисляет bounds послойно. Для $i$-го слоя
с весами $W_i$ и входным интервалом $[\ell_{i-1}, u_{i-1}]$:
\begin{align*}
\ell_i &= W_i^+ \ell_{i-1} + W_i^- u_{i-1} + b_i, \\
u_i &= W_i^+ u_{i-1} + W_i^- \ell_{i-1} + b_i,
\end{align*}
где $W_i^+ = \max(W_i, 0)$, $W_i^- = \min(W_i, 0)$.

IBP теряет информацию о корреляции между компонентами промежуточного
вектора: на каждом слое он рассматривает каждый элемент интервала как
независимую переменную. Это приводит к \textit{эффекту обёртки}
(wrapping effect)\cite{moore2009interval}~--- артефакту интервальной
арифметики, при котором ширина bounds может расти экспоненциально
с числом слоёв.

\textbf{Пример.} Рассмотрим двумерный вход $\mathbf{x} = (x_1, x_2)^T$
с $x_1, x_2 \in [-1, 1]$ и матрицу
$W = \begin{pmatrix} 1 & 1 \\ 1 & -1 \end{pmatrix}$.

CROWN: $y = W\mathbf{x}$, откуда $y_1 = x_1 + x_2 \in [-2, 2]$,
$y_2 = x_1 - x_2 \in [-2, 2]$. Повторное применение $W$:
$z_1 = y_1 + y_2 = 2x_1 \in [-2, 2]$.

IBP: первый слой даёт $y_1 \in [-2, 2]$, $y_2 \in [-2, 2]$~---
совпадает. Но при повторном применении $W$ IBP считает $y_1$ и $y_2$
\textit{независимыми}: $z_1 = y_1 + y_2 \in [-4, 4]$~---
в два раза шире истинного. CROWN, отслеживая зависимость
$y_1 = x_1 + x_2$, $y_2 = x_1 - x_2$, даёт точный результат.

\subsubsection{Количественный эффект: рост ошибки с глубиной}

Для случайной весовой матрицы с элементами порядка $O(1/\sqrt{h})$
(типичная инициализация) ширина IBP-интервала после $k$ слоёв
растёт как $O(\rho^k)$, где $\rho$~--- спектральный радиус обобщённой
матрицы абсолютных значений $|W|$.

Экспериментальные данные (Таблица~\ref{tab:tp-correctness})
подтверждают эту зависимость.
При 1 паре Col-Row (модель 256$\times$2) зона не содержит промежуточных
ReLU, IBP fallback не используется, отклонение $\sim 10^{-8}$
(машинный эпсилон \texttt{float32}).
При 2 парах (256$\times$4) каждая зона содержит промежуточный ReLU,
IBP задействован, отклонение $\sim 2.5$~--- рост на 8 порядков.
При 3 парах (256$\times$6) три зоны, каждая с ReLU, отклонение $\sim 750$.

Потеря точности обусловлена не арифметическими ошибками распределённых
вычислений, а \textit{заменой метода}: внутри шардированных зон точный
CROWN вынужденно заменяется на грубый IBP, и с каждой дополнительной
зоной ошибка накапливается мультипликативно по глубине.

TP-верификация эффективна для моделей с малой глубиной (1--2
шардированных зоны). Для глубоких моделей ($\geq 4$ зон) потеря
точности делает bounds практически бесполезными~--- что мотивирует
FSDP-подход (раздел~\ref{sec:fsdp-memory-theory}), свободный от
данного ограничения.

\subsection{Сравнение TP и FSDP для верификации}

Pipeline Parallelism для верификации неприменим: CROWN выполняет
обратный проход по \textit{всем} слоям, и при PP, где каждый GPU
хранит только свою стадию, полный обратный проход потребовал бы
пересылки весов между GPU~--- фактически превращая PP в менее
эффективный вариант TP с point-to-point коммуникацией.

FSDP шардирует веса по GPU, собирая их \texttt{AllGather}
непосредственно перед каждым слоем. В контексте верификации это
принципиальное преимущество перед TP: каждый слой работает с полной
весовой матрицей, поэтому вычисления математически \textit{идентичны}
однопроцессорному режиму и потерь точности нет.
Ценой является больший коммуникационный overhead:
\texttt{AllGather} перед каждым слоем против одного
\texttt{AllReduce} на пару слоёв в TP.

TP шардирует и веса, и A-матрицы, давая до ${\sim}2\times$ снижение
пиковой памяти, но с падением точности bounds по мере роста числа
шардированных зон.
FSDP шардирует \textit{только веса}: пиковая память снижается на
25--39\%, baseline-память (хранение весов)~--- на 80--90\%,
а bounds побитово идентичны однопроцессорному режиму, что критично
для полной верификации ($\beta$-CROWN).

\subsection{Реализация TP для bound propagation в \texttt{auto\_LiRPA}}

Для интеграции TP в \texttt{auto\_LiRPA}\cite{AutoLiRPARepo}
реализованы четыре компонента.

Операторы \texttt{BoundLinearTP\_Col} и \texttt{BoundLinearTP\_Row}
наследуют от \texttt{BoundLinear}, хранят шардированные весовые
матрицы и добавляют \texttt{AllReduce} в \texttt{bound\_backward},
обеспечивая корректное распространение границ через распределённые слои.

Каждый оператор регистрируется как пользовательская операция ONNX
через \texttt{torch.autograd.Function} и метод \texttt{symbolic},
что позволяет JIT-трассировке PyTorch корректно строить вычислительный
граф.

Функция \texttt{tp\_shard\_bounded\_module} принимает произвольный
\texttt{BoundedModule} (включая загруженный из ONNX) и автоматически
заменяет чередующиеся пары \texttt{BoundLinear} на
\texttt{BoundLinearTP\_Col/Row}, шардируя весовые матрицы по числу GPU.

Между каждой парой Column--Row узлов формируется <<шардированная
зона>>~--- подграф с уменьшенной размерностью активаций.
Узлы в зонах (включая ReLU) обрабатываются локально без
синхронизации между GPU; определение принадлежности узла зоне
выполняется BFS-разметкой при инициализации модуля.

\subsection{Реализация FSDP для bound propagation в \texttt{auto\_LiRPA}}

FSDP-реализация не требует новых типов операторов и не изменяет
вычислительный граф~--- весь объём составляет около 100 строк кода
против ${\sim}500$ для TP.

Функция \texttt{fsdp\_shard\_bounded\_module} обходит граф
\texttt{BoundedModule}, находит узлы \texttt{BoundParams}
полносвязных (\texttt{BoundLinear}) и свёрточных (\texttt{BoundConv})
слоёв и заменяет каждую матрицу на шард по нулевому измерению
(\texttt{dim=0}).
Для \texttt{BoundLinear} матрица $W \in \mathbb{R}^{k \times n}$
заменяется на $W^{(r)} \in \mathbb{R}^{(k/P) \times n}$;
для \texttt{BoundConv} матрица
$W \in \mathbb{R}^{C_\text{out} \times C_\text{in} \times k_H \times k_W}$
заменяется на $W^{(r)} \in \mathbb{R}^{(C_\text{out}/P) \times C_\text{in} \times k_H \times k_W}$.
Шардирование по выходному измерению корректно, поскольку оба типа
слоёв линейны по выходу.

При вызове метода \texttt{forward()} класса \texttt{BoundParams}
обнаружение флага шардирования инициирует \texttt{dist.all\_gather},
собирающий полную матрицу из шардов всех GPU.
Полная матрица возвращается как результат forward и используется
всеми нижестоящими операциями без модификации.

Хуки \texttt{fsdp\_gather\_node} и \texttt{fsdp\_free\_node}
интегрированы в \texttt{backward\_general} (CROWN backward) и
\texttt{IBP\_general} (IBP forward): полная матрица собирается
непосредственно перед обработкой слоя и освобождается сразу после.
В каждый момент времени в GPU-памяти находится не более одной полной
весовой матрицы (помимо шардов).

\subsubsection{Теоретический анализ расхода памяти}
\label{sec:fsdp-memory-theory}

Пусть модель содержит $d$ линейных слоёв с весовыми матрицами
$W_i \in \mathbb{R}^{h \times h}$, верификация выполняется методом
CROWN для одного примера, $P$~--- число GPU.

В single-GPU режиме все веса хранятся целиком: $M_W = d \cdot h^2$.
При FSDP каждый GPU хранит $1/P$ каждого веса плюс одну полную
матрицу для текущего слоя:
$$M_W^{\text{FSDP}} = \frac{d \cdot h^2}{P} + h^2.$$

В CROWN backward для промежуточного слоя $i$ хранится A-матрица
размера $O(h \times h)$. При TP A-матрицы проходят через
шардированные слои с уменьшенной размерностью $h/P$; при FSDP они
всегда имеют полную размерность~$h$:
$$M_A^{\text{TP}} \propto \frac{h^2}{P}, \qquad
  M_A^{\text{FSDP}} \propto h^2.$$

Суммарная пиковая память для TP:
$$M_{\text{peak}}^{\text{TP}} \approx
    \frac{(d + C) \cdot h^2}{P},$$
что даёт линейное снижение в $P$ раз. Для FSDP:
$$M_{\text{peak}}^{\text{FSDP}} \approx
    h^2 \!\left(\frac{d}{P} + 1 + C\right),$$
где $C$~--- число одновременно хранящихся A-матриц.

FSDP экономит память только на весах, тогда как TP сокращает оба
слагаемых. Отношение пиковой памяти FSDP к single-GPU:
$$\frac{M_{\text{peak}}^{\text{FSDP}}}{M_{\text{peak}}^{\text{single}}}
    = \frac{d/P + 1 + C}{d + C}.$$
При $d = 4$, $C \approx 4$, $P = 2$ получаем
$\approx (2 + 1 + 4)/(4 + 4) = 7/8 = 0.875$~---
что соответствует экспериментально наблюдаемой экономии $\sim$25--34\%.

Baseline-память (хранение весов без A-матриц) сокращается сильнее:
$$\frac{M_W^{\text{FSDP}}}{M_W^{\text{single}}}
    = \frac{1}{P} + \frac{1}{d},$$
что при $d = 4$, $P = 2$ даёт $0.75$, а при $d = 8$, $P = 2$~---
$0.625$, соответствуя экспериментальным $\sim$82--90\%.

\subsection{Трудности интеграции TP с полной верификацией}

BaB с TP представляет значительно более сложную задачу, чем с FSDP.
Во-первых, в Branch-and-Bound каждый шаг выбирает нейрон для фиксации
состояния. При TP нейроны шардированных слоёв распределены по GPU,
и выбор нейрона для ветвления требует глобальной координации:
сбора статистик со всех GPU и рассылки решения.
Во-вторых, $\beta$-CROWN кодирует ограничения ветвления через
двойственные переменные $\beta$; для нейронов в шардированных зонах
$\beta$ определены локально, но их влияние распространяется через
\texttt{AllReduce} на глобальные bounds, что требует аккуратной
координации обновлений.
В-третьих, BaB генерирует миллионы подзадач; каждая из них в батче
должна использовать одинаковую TP-конфигурацию, что усложняет
балансировку нагрузки.

FSDP свободен от перечисленных трудностей: каждый ранк получает
полные весовые матрицы через AllGather и выполняет логику ветвления,
идентичную однопроцессорной.
Интеграция FSDP с $\beta$-CROWN + BaB реализована и экспериментально
подтверждена (см. раздел~\ref{sec:exp6}): при идентичном workload
оба режима~--- single-GPU и FSDP$=$2~--- дают побитово совпадающие
результаты с overhead AllGather не более 5~МБ.
